% =====================================================================
%  CAPÍTULO 7 — CONCLUSÃO
% =====================================================================

\chapter{CONCLUSÃO}
\label{cap:conclusao}

Este trabalho apresentou o registro de software do \argus{}, um simulador
multiplataforma de ciclismo virtual de código aberto, e documentou a evolução
do projeto desde a versão inicial registrada junto ao INPI até a versão atual.
O desenvolvimento demonstrou a viabilidade técnica de uma plataforma de
ciclismo virtual gratuita, multiplataforma e orientada ao processamento local.

\section{Considerações finais}
\label{sec:consideracoes-finais}

O \argus{} diferencia-se por combinar simulação física em tempo real, análise
fisiológica avançada, integração com dispositivos BLE e recursos de
inteligência artificial executados localmente, preservando a privacidade dos
usuários e reduzindo dependências de infraestrutura em nuvem. Adicionalmente, o
sistema de rastreamento de desgaste de componentes oferece uma ferramenta
prática de manutenção preventiva, preenchendo uma lacuna não endereçada pelas
plataformas comerciais existentes.

A arquitetura em \textit{Clean Architecture}, com os padrões \textit{Facade} e
\textit{Repository}, demonstrou-se adequada para o domínio: permitiu a evolução
contínua do sistema --- como a adição do aplicativo Android, das integrações
com o Strava e o Intervals.icu e do assistente de IA --- sem comprometer a
estabilidade dos módulos existentes. A validação experimental com equipamentos
reais confirmou o atendimento a todos os requisitos funcionais e não
funcionais definidos.

Além da contribuição tecnológica, a proposta possui relevância social
significativa ao democratizar o acesso ao treinamento esportivo científico,
especialmente para atletas amadores e usuários com limitações financeiras ou
computacionais. Dessa forma, o \argus{} consolida-se como uma alternativa
inovadora às soluções proprietárias do mercado, unindo acessibilidade,
desempenho e liberdade tecnológica.

\section{Rastreabilidade de objetivos}
\label{sec:rastreabilidade}

A Tabela~\ref{tab:rastreabilidade} apresenta a relação entre os objetivos
específicos definidos no Capítulo~\ref{cap:introducao}, as implementações
realizadas e as evidências de validação.

\begin{table}[h]
    \centering
    \caption{Rastreabilidade: objetivos, implementação e validação}
    \label{tab:rastreabilidade}
    \small
    \begin{tabular}{
        >{\raggedright\arraybackslash}p{3.2cm}
        >{\raggedright\arraybackslash}p{5cm}
        >{\raggedright\arraybackslash}p{5cm}
    }
        \toprule
        \textbf{Objetivo} & \textbf{Implementação} &
        \textbf{Validação} \\
        \midrule
        Conectividade BLE & BLE Service
        (\code{internal/service/ble}) & Tabela~\ref{tab:validacao}
        (estabilidade) \\
        Motor de simulação & Equação de Martin
        (\code{internal/service/sim}) & 200+ arquivos \texttt{.fit}
        aceitos \\
        Interface 60fps & \code{TelemetryEventBus} &
        Teste manual (HUD fluido) \\
        Métricas fisiológicas & NP, TSS, PMC, TRIMP
        (\code{internal/service/fit}) &
        Figura~\ref{fig:screenshot-career} \\
        IA Coach local & Ollama
        (\code{internal/service/ai}) &
        Figura~\ref{fig:screenshot-desgaste} \\
        Desgaste componentes & \code{BikeComponent} &
        Figura~\ref{fig:screenshot-desgaste} \\
        SQLite multiperfil & \code{StorageFacade} &
        Tabela~\ref{tab:validacao} \\
        Exportação \texttt{.fit} & \code{muktihari/fit} &
        Strava e Intervals.icu \\
        Multiplataforma & Wails + Capacitor &
        Windows, Linux, Android \\
        Gamificação & \code{gamification.go} &
        Figura~\ref{fig:screenshot-trophy} \\
        \bottomrule
    \end{tabular}
\end{table}

\section{Contribuições do trabalho}
\label{sec:contribuicoes}

As principais contribuições deste trabalho são:

\begin{itemize}
    \item o desenvolvimento de um simulador de ciclismo virtual multiplataforma,
        gratuito e de código aberto (GPLv3), com processamento local
        \textit{offline-first};
    \item a implementação de um motor de simulação física baseado na equação de
        Martin, com controle bidirecional de rolos inteligentes via FTMS e FE-C;
    \item a implementação local de métricas fisiológicas avançadas (NP, IF,
        TSS, TRIMP, PMC e desvio cardiovascular) normalmente restritas a
        plataformas pagas;
    \item o assistente de treino baseado em IA generativa local, com geração de
        planos ZWO e preservação de privacidade;
    \item o sistema de rastreamento de desgaste de componentes mecânicos da
        bicicleta;
    \item o registro do software junto ao INPI (Processo BR512026003786-7),
        formalizando a proteção intelectual da versão inicial do sistema.
\end{itemize}

\section{Trabalhos futuros}
\label{sec:trabalhos-futuros}

Como perspectivas de evolução do projeto, destacam-se:

\begin{itemize}
    \item \textbf{Suporte multiplataforma completo em dispositivos móveis}:
        ampliar a validação do aplicativo Android e avaliar a compilação para
        iOS;
    \item \textbf{Modo \textit{multiplayer} local}: permitir que atletas
        treinem juntos na mesma rede local, sem infraestrutura em nuvem;
    \item \textbf{Mapas totalmente \textit{offline}}: aprimorar o
        \textit{caching} de \textit{tiles} e disponibilizar pacotes de mapas
        para download;
    \item \textbf{Expansão do modelo de desgaste}: incorporar variáveis como
        tipo de treino e intensidade ao cálculo de desgaste dos componentes;
    \item \textbf{Melhorias no assistente de IA}: uso de \textit{retrieval
        augmented generation} sobre o histórico do atleta e otimização da
        inferência para dispositivos sem GPU;
    \item \textbf{Ampliação da validação experimental}: testes com uma maior
        diversidade de fabricantes de rolos e sensores, bem como a publicação
        de um estudo de usabilidade com atletas amadores.
\end{itemize}
